\documentclass[11pt,letterpaper]{article}
\usepackage[margin=1in]{geometry}
\usepackage[T1]{fontenc}
\usepackage{lmodern,microtype}
\usepackage{amsmath,amssymb,amsthm,mathtools}
\usepackage{booktabs,longtable,array}
\usepackage{xcolor,fancyhdr,xurl}
\usepackage[hidelinks]{hyperref}
\usepackage{listings}
\definecolor{inkblue}{HTML}{163D5B}
\definecolor{muted}{HTML}{4F6070}
\hypersetup{pdftitle={Coprime product certificates and positive Pell centers for Diophantine consistency},pdfauthor={Research prepared for Vladimir Reshetnikov with OpenAI assistance}}
\pagestyle{fancy}\fancyhf{}
\fancyhead[L]{\small\color{muted}Diophantine consistency: new explicit reductions}
\fancyhead[R]{\small\color{muted}25 October 2026}
\fancyfoot[C]{\thepage}
\setlength{\headheight}{14pt}
\setlength{\parindent}{0pt}\setlength{\parskip}{5pt}
\setlength{\emergencystretch}{2em}
\allowdisplaybreaks[2]
\newtheorem{theorem}{Theorem}[section]
\newtheorem{lemma}[theorem]{Lemma}
\newtheorem{proposition}[theorem]{Proposition}
\newtheorem{corollary}[theorem]{Corollary}
\theoremstyle{remark}\newtheorem{remark}[theorem]{Remark}
\newcommand{\N}{\mathbb N}
\newcommand{\Z}{\mathbb Z}
\newcommand{\Pos}{\mathbb Z_{>0}}
\newcommand{\Con}{\operatorname{Con}}
\newcommand{\ZFC}{\mathsf{ZFC}}
\newcommand{\fl}[1]{\left\lfloor#1\right\rfloor}
\newcommand{\rg}[2]{\{#1,\ldots,#2\}}
\newcommand{\file}[1]{\texttt{\small\detokenize{#1}}}
\lstset{basicstyle=\ttfamily\small,breaklines=true,columns=fullflexible,frame=single,rulecolor=\color{black!15}}
\title{\color{inkblue}\textbf{Coprime product certificates and positive Pell centers}\\[5pt]
\large Explicit improvements for Diophantine formulations of $\Con(\ZFC)$}
\author{Research prepared for Vladimir Reshetnikov and ProveIt\\
\small with OpenAI assistance}
\date{25 October 2026}

\begin{document}
\maketitle
\begin{abstract}
We construct several fixed ordinary integer polynomials whose lack of solutions,
over their explicitly stated domains, is provably equivalent in ZFC to
syntactic consistency of ZFC. The central reduction replaces separate square
witnesses by a single product certificate. A prescribed divisibility condition
on a freely chosen Pell witness creates the coprimality needed for exact prime
valuation recovery. Applied to the existing fifteen-variable, 66-operation
construction, staged square recovery gives fourteen positive variables in 70
operations at degree 189, and a 71-operation variant at degree 157. Two
divisibility anchors give thirteen positive variables in 72 operations at
degree 433.

A separate ten-positive-variable construction reuses $XY$ as its congruence
carrier. A bounded Pell index is recovered before full denominator divisibility
is asserted. Together with a proved unshifted quotient, this gives 76 operations
at degree 626, or 77 at degree 464; delaying one divisibility condition gives
80 operations at degree 382. Finally, a new positive polynomial comparison
center, pooled digit bounds, and a noncircular bootstrap deletion produce a
direct C222 quartic in 221 nonnegative variables, with 172 quadratic residuals
and 9,536 collected monomials. This last route uses neither K215 nor Jones's
universal polynomial.

The archive supplies literal equations, all fixed coefficient data, ordinary
proofs, exact reconstruction programs, and independently checked finite
certificates. These are bounded improvements relative to the pinned repository
and incoming archives. They are not claims of global simplicity, a proof of
consistency, or a new end-to-end proof-assistant derivation.
\end{abstract}
\thispagestyle{empty}
\clearpage
\tableofcontents
\clearpage

\section{What is constructed and what is proved}

Fix the recursive presentation of ZFC used by the retained ProveIt calibration.
The arithmetic sentence $\Con(\ZFC)$ means that there is no finite syntactic
refutation in that presentation. The objective is a fixed polynomial over
$\Z$, together with a specified witness domain, and a proof of
\[
\ZFC\vdash
\bigl(\neg\exists\boldsymbol{x}\in D^n\;P(\boldsymbol{x})=0\bigr)
\ \longleftrightarrow\ \Con(\ZFC).
\]
All new reductions in this report are ordinary mathematical proofs formalizable
in ZFC. Their final logical calibration is inherited from named, retained
ProveIt source results. The finite checkers verify the specified polynomials
and their algebraic transformations. They do not decide whether a refutation
exists, and no such refutation is assumed.

\subsection{Results and comparisons}

The comparison snapshot is the full commit
\begin{center}\small\texttt{e0c98c629bf10d30b695884f6b1d2ab361bdf513}.\end{center}
It includes the active research, the manuscript collection, and all six ZIP
archives in \file{docs/incoming}. The collection overview alone is older than
several actual endpoints. In particular, the incoming certificate-reduction
archive already supplies ten-positive 78/630 and 79/468 circuits, and the
incoming graph-digit archive supplies a 226-variable direct natural quartic.
These stronger objects are included in our comparisons \cite{repo,incoming}.

\begin{center}\small
\begin{tabular}{@{}lrrrr@{}}
\toprule
Construction & Domain & Variables & Operations & Degree\\
\midrule
Coprime product, staged recovery & $\Pos$ & 14 & \textbf{70} & \textbf{189}\\
Coprime product, reduced coefficient & $\Pos$ & 14 & \textbf{71} & \textbf{157}\\
Product certificate with two anchors & $\Pos$ & 13 & \textbf{72} & \textbf{433}\\
$XY$ carrier, square transport & $\Pos$ & 10 & \textbf{76} & \textbf{626}\\
$XY$ carrier, bilinear transport & $\Pos$ & 10 & \textbf{77} & \textbf{464}\\
Delayed divisibility, bilinear transport & $\Pos$ & 10 & \textbf{80} & \textbf{382}\\
Direct positive-center quartic & $\N$ & \textbf{221} & --- & \textbf{4}\\
Sparse direct companion & $\N$ & \textbf{222} & --- & \textbf{4}\\
\bottomrule
\end{tabular}
\end{center}

Here $\N=\{0,1,2,\ldots\}$, whereas $\Pos$ excludes zero. A named arithmetic
register is an abbreviation for a polynomial, not an additional unknown.
Each counted operation is one live binary addition, subtraction, or
multiplication; squares and fixed scalar multiplications are charged. These
are costs of supplied evaluators, not lower bounds on arithmetic complexity.

The fourteen-variable 70/189 form improves both cost and degree over the
inspected 73/342 point; the 71/157 form improves both over 76/191. The
thirteen-variable 72/433 form lowers the inspected minimum operation count of
83, while accepting a higher degree than the 83/114 form. The ten-variable
76/626 and 77/464 forms improve both coordinates of the incoming 78/630 and
79/468 forms. Existing nine-variable equations remain smaller in number of
unknowns, with substantially larger evaluators and degrees. The maintained
fifteen-variable 66/157 equation remains cheaper in operations than our new
thirteen- and fourteen-variable forms. Thus these results have distinct uses;
there is no single metric-independent winner.

\begin{center}\small
\begin{tabular}{@{}lrrrr@{}}
\toprule
Direct natural quartic & Variables & Residuals & Terms & Height\\
\midrule
Inspected graph-digit parent & 226 & 175 & 10,314 & 87,839,583\\
New positive-center form & \textbf{221} & \textbf{172} & \textbf{9,536} & \textbf{83,954,186}\\
Inspected sparse parent & 227 & 176 & 7,905 & 87,839,583\\
New sparse companion & \textbf{222} & \textbf{173} & \textbf{6,821} & \textbf{83,954,028}\\
\bottomrule
\end{tabular}
\end{center}
Height here is the largest absolute coefficient after collecting the quartic.
These comparisons keep both the domain and degree fixed. They concern the
direct C222 route, not the smallest known quartic under every possible encoding.

\subsection{Novelty and inherited material}

The new contribution is a set of proved constructions and their exact
compositions. Prime valuation arguments, Pell recurrences, product-square
recovery, and sums of squares are classical tools. The source proof calculus,
native compiler, fifteen-variable source, doubled strong pair, semantic anchor,
and self-sized radix have explicit predecessors in ProveIt. We attribute them
accordingly. The newly imposed cofinal multiplier, its low-cost square-witness
elimination, the $XY$ carrier with its signed unshifted proof, and the positive
comparison center are the additional mechanisms supplied here. No claim of
literature-wide priority follows from the repository audit.

Universal Diophantine constructions explain why such equations exist, but do
not by themselves instantiate the small fixed equations required here
\cite{jones,bayer}. Every parameter needed for our actual polynomials is fixed
by literal data or by the complete finite recipe in Appendix~\ref{app:constants}.

\section{Two elementary certificate lemmas}
\label{sec:lemmas}

\begin{lemma}[Coprime product recovery]\label{lem:product}
Let $G,Q,V,t$ be integers with $G>0$, $Q>1$, $t>0$, $V\ne0$, and
\[
\gcd\bigl(G,(Q-1)(QV^2-1)\bigr)=1.
\]
If
\begin{equation}\label{eq:product}
(Q-1)t^2=G(QV^2-1),
\end{equation}
then $G=g^2$ for a positive integer $g$, $g\mid t$, and $y=t/g>0$ is an
integer satisfying
\[
QV^2-(Q-1)y^2=1.
\]
Conversely, these conclusions imply \eqref{eq:product} with $t=gy$.
\end{lemma}
\begin{proof}
For every prime $p\mid G$, the coprimality assumption removes both other
factors from the valuation equation, leaving $v_p(G)=2v_p(t)$. Thus every
prime exponent in $G$ is even, and the positive square root $g$ divides $t$.
Divide \eqref{eq:product} by $g^2$ and rearrange. The converse is substitution.
\end{proof}

It is not enough to declare a product square: nonsquare factors may share a
squareclass. The construction must provide the displayed coprimality. Here
this is achieved economically either by imposing divisibility on a freely
chosen Pell ordinate or by making $Q$ itself a multiple of $G$.

The useful expanded residual is
\begin{equation}\label{eq:folded}
\mathcal E_G=Q(t^2-GV^2)-t^2+G.
\end{equation}
It is identically $(Q-1)t^2-G(QV^2-1)$. This form shares a multiplication
and avoids separately computing both subtracted ones.

\begin{lemma}[A unit isolates a residual]\label{lem:unit}
If $e,z_1,\ldots,z_r$ are integers and
\[
(3e+1)z_1\cdots z_r=1,
\]
then $e=0$ and $z_1\cdots z_r=1$.
\end{lemma}
\begin{proof}
Every factor is $1$ or $-1$. The value $3e+1=-1$ would give $3e=-2$, which
is impossible over the integers. Therefore $3e+1=1$.
\end{proof}

\begin{lemma}[Cofinal divisible Pell ordinates]\label{lem:period}
For $A\ge2$ and $m\ge1$, there are arbitrarily large positive indices $n$
such that $m\mid Y_n(A)$, where
\[
X_0=1,\quad Y_0=0,\qquad
X_{n+1}=AX_n+(A^2-1)Y_n,\quad Y_{n+1}=X_n+AY_n.
\]
\end{lemma}
\begin{proof}
The recurrence matrix has determinant one. Modulo $m$ it is an invertible
element of a finite matrix group, so a positive power is the identity.
Every positive multiple of that period sends $(1,0)$ to $(1,0)$ modulo $m$.
The ordinates grow strictly for positive indices.
\end{proof}
This is why adding a divisibility condition on a private Pell witness need
not restrict which outer computations or refutations can be represented.
The completeness proof must choose fresh auxiliary witnesses; it must not
assume arbitrary old witnesses already satisfy the new condition.

\section{The common fifteen-variable source}
\label{sec:source}

This section records the precise arithmetic source used by the new
thirteen- and fourteen-variable constructions. Its fixed ports
$\beta,K_0,M_C,M_F$ are the entire indexed943/radix67 tuple of the retained
native compiler. They are not quantified variables. Appendix~\ref{app:constants}
prints their finite recipe; $K_0$ is denoted $K$ in that recipe, and $M_F$ is
the already shifted source port. Ports from different native tuples must not
be interchanged.

For positive outer coordinates define
\begin{align*}
q_-&=\beta J,& q&=q_-+1,& C&=q-F-\alpha,\\
X&=wq,& Y&=sq^3,& N&=XY,\\
k&=\eta+\zeta,& c&=kY+\eta,& \ell&=NkY,\\
\Omega&=1+\ell(\ell+k),& a&=N+Y,& A&=a+2,\\
H&=4a+3,& \Delta&=a^2+H,& D&=X+ac+\gamma H,\\
\end{align*}
\begin{align*}
R&=[q_-(q-F)+\alpha](q^2-1)+(M_C+qM_F)J,\\
N_1&=D^2-\Delta c^2,\\
N_k&=k-hN-R,\\
N_t&=(K_0+w)C+(q-F)-zq_-.
\end{align*}
Thus $\Delta=A^2-1$. In particular $A>2$, $\Delta>0$, $c>0$, and $\Omega>1$ before imposing
any equation. The original fifteen positive variables are
\[
J,F,\alpha,z,f,h,i_0,T,s,w,\tau,\eta,\zeta,y,\gamma.
\]
With $S=i_0\Delta c^2$ and $V=f^2(cT-R)-c$, its six factors are
\begin{align*}
N_0&=\tau^2-\ell(\ell+k),&
N_a&=S^2(V^2-y^2)+y^2,\\
N_s&=\Delta f^2-S^2,&
P_{66}&=N_0N_1N_aN_kN_tN_s-\Delta.
\end{align*}
This is the accepted fifteen-positive, 66-operation, degree-157 free-port
polynomial, specialized to the complete indexed tuple \cite{p66}.

\subsection{The inherited theorem used here}

The retained ordinary source theorem provides both directions:
\begin{enumerate}
\item Every positive zero of $P_{66}$ decodes a finite syntactic ZFC refutation.
\item A finite refutation supplies canonical positive outer coordinates with
$\tau^2=\Omega$, $D^2-\Delta c^2=1$, $N_k=N_t=1$,
$c=Y_R(A)$, and $R\equiv3\pmod4$. The private strong Pell pair and auxiliary
quotient/ordinate can be reconstructed as follows.
\end{enumerate}
Choose any positive pair $f,I$ with $f^2-\Delta I^2=1$ and $c^2\mid I$.
Put $i_0=I/c^2$ and $S=\Delta I$. There are positive integers $V,y,T$ with
\begin{equation}\label{eq:aux-source}
S^2(V^2-y^2)+y^2=1,\qquad V=f^2(cT-R)-c.
\end{equation}
For explicit reconstruction take
\[
V=\frac{X_R(S)}S,\qquad y=Y_R(S),\qquad
T=\frac{V+c+Rf^2}{cf^2}.
\]
These divisions occur in a proof of existence, not in a polynomial definition.
The retained squared-modulus lemma proves their integrality. Briefly, the
odd-index Chebyshev polynomials $Q_j$ satisfy
\[
X_{2j+1}(s)=sQ_j(s^2),\quad Q_j(0)=(-1)^j(2j+1),\quad
Q_j(1-A^2)=(-1)^jY_{2j+1}(A).
\]
Since $S^2=\Delta(f^2-1)$ and $R\equiv3\pmod4$, this gives
$V\equiv-c\pmod{f^2}$ and $V\equiv-R\pmod c$.
Also $\gcd(c,f)=1$ and $f^2\equiv1\pmod c$, proving the displayed formula
for $T$. This uniform reconstruction is the precise completeness interface
needed below. Its ordinary proof and the native specialization are retained
in the artifact sources; no new formal-kernel calibration is asserted.

\section{Fourteen variables by staged square recovery}
\label{sec:fourteen}

The unknowns are
\[
J,F,\alpha,z,f,h,i,T,s,w,\eta,\zeta,\gamma,t\in\Pos.
\]
Keep the common definitions in Section~\ref{sec:source}, and put
$I=ic^2$ and $V=f^2(cT-R)-c$. The two complete choices are
\[
\begin{array}{c|cc}
 & W & Q\\\hline
\text{staged form} & \Delta\Omega I^2 & \Delta W\\
\text{reduced form} & \Delta(\Omega I)^2 & \Delta(f^2-1)
\end{array}
\qquad L_s=f^2-W.
\]
For either choice define
\[
\mathcal E=Q(t^2-\Omega V^2)-t^2+\Omega,
\qquad
\boxed{P_{14}=(3\mathcal E+1)N_1N_kN_tL_s-1=0.}
\]
This is an ordinary integer polynomial. The staged version uses 70 operations
and has degree 189; the reduced version uses 71 operations and has degree 157.
The versions have different private witness parametrizations. They are
existence-equivalent, not identical polynomials or asserted to have the same
auxiliary-coordinate zero sets.

\begin{theorem}\label{thm:fourteen}
For either displayed $P_{14}$, existence of a positive zero is equivalent to
existence of a positive zero of the calibrated source, and hence to a finite
ZFC refutation. Consequently ZFC proves the requested no-zero equivalence
with $\Con(\ZFC)$.
\end{theorem}
\begin{proof}[Soundness]
Every factor in the unit product is $\pm1$.
Lemma~\ref{lem:unit} gives $\mathcal E=0$.
Because $\Delta=A^2-1$ is zero or three modulo four, $N_1$ cannot be $-1$,
so $N_1=1$. The order for the remaining factors depends on the version.

In the staged form, $Q=\Delta^2\Omega I^2>1$ and $\Omega\mid Q$ before
knowing the sign of $L_s$. The equation $\mathcal E=0$ itself excludes
$V=0$, since its left side would be $(Q-1)t^2+\Omega>0$.
Lemma~\ref{lem:product} first yields $\Omega=\tau^2$ and $y=t/\tau>0$.
Only now write $L_s=f^2-\Delta(\tau I)^2$; reduction modulo four excludes
$L_s=-1$, so $L_s=1$ and $N_kN_t=1$.
Set $i_0=\tau i$. Then $Q=(i_0\Delta c^2)^2$ and the recovered auxiliary
norm is exactly the source norm.

In the reduced form, $L_s=f^2-\Delta(\Omega I)^2$ already excludes $-1$.
Thus $L_s=1$ first, and $Q=\Delta^2(\Omega I)^2>1$ with $\Omega\mid Q$.
Again $\mathcal E=0$ excludes $V=0$, and the same lemma supplies
$\tau=\sqrt\Omega$ and $y=t/\tau$. Here set $i_0=\Omega i$.
In both versions $Q(V^2-y^2)+y^2=1$.

The restored source factors satisfy $N_0=N_1=N_a=1$,
$N_s=\Delta$, and $N_kN_t=1$. Thus the source polynomial vanishes.
Every restored coordinate is strictly positive.
\end{proof}
\begin{proof}[Completeness]
Start with canonical outer witnesses supplied by a finite refutation,
including $\Omega=\tau^2$ and $c=Y_R(A)$ with $R\equiv3\pmod4$.
By Lemma~\ref{lem:period}, choose a positive strong Pell pair $f,I_0$.
For the staged version require $\tau c^2\mid I_0$ and put
$i=I_0/(\tau c^2)>0$; for the reduced version require
$\Omega c^2\mid I_0$ and put $i=I_0/(\Omega c^2)>0$.
The uniform reconstruction \eqref{eq:aux-source} supplies new positive
$T,V,y$; let $t=\tau y$. Now $L_s=N_1=N_k=N_t=1$, and
\eqref{eq:product} gives $\mathcal E=0$ in both versions.
This does not require an arbitrary old multiplier to have satisfied the
new divisibility condition.
\end{proof}

\subsection{The useful dependency order}

The staged form places only one copy of $\Omega$ in $W$.
It would be invalid to call $W/\Delta$ a square before recovering
$\Omega=\tau^2$. Instead, the product certificate first makes $\Omega$
a square, and the strong unit sign is resolved afterward. This lowers the
degree from the more direct squared-multiplier construction while preserving
its operation count. The proof depends on this order, not on an informal
replacement of a factor by its anticipated square root.

\subsection{Exact degree}

Assign degree one to the unknowns and zero to the fixed ports. The relevant
exact degrees are
\[
\deg\Omega=22,\quad\deg c=5,\quad\deg\Delta=12,\quad\deg V=8.
\]
The main norm has degree 18 rather than its naive upper bound 22. With
$Z=X+\gamma H$, the exact identity is
\[
D^2-\Delta c^2=2acZ+Z^2-Hc^2,
\]
and the leading term comes from $2ac\gamma H$.
Also $\deg N_k=7$ and $\deg N_t=2$.
For the staged form, $\deg Q=68$, $\deg\mathcal E=106$, and
$\deg L_s=56$, giving $106+18+7+2+56=189$.
For the reduced form, $\deg Q=14$, $\deg\mathcal E=52$, and
$\deg L_s=78$, giving 157.
The leading-term certificate is nonzero at the authentic positive $\beta$;
the supplied exact checks also verify the literal circuit and its degree.


\section{Thirteen variables with two divisibility anchors}
\label{sec:thirteen}

This construction removes $f$ as an independent unknown as well as the two
separate square witnesses. Its variables are
\[
J,F,\alpha,z,h,i,T,s,w,\eta,\zeta,\gamma,t\in\Pos.
\]
Retain the common prefix of Section~\ref{sec:source}. In particular,
$A=a+2$ is the main Pell base and $D=X+ac+\gamma H$ its retained
candidate abscissa. Define
\begin{align*}
j&=ADi,& U&=jc^2,& W&=\Delta\Omega U^2,\\
F_2&=1+W,& Q&=\Delta W,& G&=\Omega F_2,\\
V&=F_2(cT-R)-c,&
\mathcal E&=Q(t^2-GV^2)-t^2+G.
\end{align*}
The equation is
\begin{equation}\label{eq:thirteen}
\boxed{P_{13}=(3\mathcal E+1)N_1N_kN_t-1=0.}
\end{equation}
Its supplied evaluator has 72 operations: 41 multiplications, 19 additions,
and 12 subtractions. Its exact total degree is 433.

\begin{theorem}\label{thm:thirteen}
The positive existence statement for \eqref{eq:thirteen} is equivalent to a
finite ZFC refutation, with the same entire fixed coefficient tuple. Therefore
\[
\ZFC\vdash
\bigl(\neg\exists\boldsymbol{x}\in\Pos^{13}\;P_{13}(\boldsymbol{x})=0\bigr)
\longleftrightarrow\Con(\ZFC).
\]
\end{theorem}
\begin{proof}[Soundness]
The unit argument gives $\mathcal E=0$, $N_1=1$, and $N_kN_t=1$.
All of $G,F_2,Q$ are positive and $Q>1$.
The equation $\mathcal E=0$ excludes $V=0$.
Put $B=Q-1$ and $C_a=QV^2-1$. The required coprimalities are proved
before recovering any missing square root.

First, $\Omega\mid Q$ and $F_2\equiv1\pmod\Omega$, so
\[
\gcd(\Omega,F_2BC_a)=1.
\]
Next $W=F_2-1$ and $V\equiv-c\pmod{F_2}$ imply
\[
B\equiv-\Delta-1=-A^2\pmod{F_2},
\]
\[
C_a\equiv-\Delta c^2-1=-D^2\pmod{F_2}.
\]
The second equality uses the already established $N_1=1$.
The two factors in $j=ADi$ ensure $F_2\equiv1\pmod A$ and
$F_2\equiv1\pmod D$. Hence $\gcd(F_2,BC_a)=1$.
Together these statements give
\[
\gcd(G,BC_a)=1,\qquad \gcd(\Omega,F_2)=1.
\]
Lemma~\ref{lem:product} supplies $G=g^2$, $g\mid t$, and
$y=t/g>0$ with $QV^2-(Q-1)y^2=1$.
The positive coprime factors of $G=\Omega F_2$ are individually squares:
\[
\Omega=\tau^2,\qquad F_2=f^2,\qquad g=\tau f.
\]
Set $I_0=\tau U$ and $i_0=I_0/c^2=\tau ADi>0$.
Then
\[
f^2-\Delta I_0^2=1,\qquad Q=(\Delta I_0)^2,
\qquad V=f^2(cT-R)-c.
\]
These restore the complete first, main, strong, and auxiliary source norms.
Together with $N_kN_t=1$, they give a positive zero of $P_{66}$.
\end{proof}
\begin{proof}[Completeness]
Begin with canonical outer data and $\tau^2=\Omega$.
By Lemma~\ref{lem:period}, choose a positive strong Pell pair $f,I_0$
such that $\tau ADc^2\mid I_0$. Put
\[
i=I_0/(\tau ADc^2),\quad U=I_0/\tau.
\]
Then $W=\Delta I_0^2=f^2-1$, $F_2=f^2$, $G=(\tau f)^2$,
and $Q=(\Delta I_0)^2$.
Reconstruct positive $T,V,y$ by \eqref{eq:aux-source}, and set
$t=\tau f y$. All factors in \eqref{eq:thirteen} have the required values.
This extra divisibility is imposed only on the freely chosen strong Pell
ordinate; the canonical outer computation is unchanged.
\end{proof}

\subsection{Why both anchors are present}

The factor $A$ controls the possible common primes of $F_2$ with $Q-1$.
The factor $D$ controls the possible common primes of $F_2$ with $QV^2-1$.
Neither coprimality is inferred merely from $F_2=1+W$.
The main norm must be recovered before using the residue $-D^2$.
This explicit ordering avoids the shared-squareclass obstruction that would
invalidate an unrestricted product-of-radicands shortcut.

The construction also avoids forcing $G\mid Q$: here only $\Omega\mid Q$.
The additional coprimalities supply exactly what the valuation argument
needs. A doubled strong-pair construction supplied an intermediate 72/465
form; the anchored version improves its degree at the same operation count.
That intermediate is not counted as another preferred endpoint.

\subsection{Exact degree and a nonzero coefficient}

The relevant register degrees are
\[
\deg j=18,\quad\deg U=28,\quad\deg W=90,\quad\deg Q=102,
\quad\deg G=112,\quad\deg V=96.
\]
Thus $-QGV^2$ is the unique highest-degree summand of $\mathcal E$, of
degree 406. Adding the factor degrees 18, 7, and 2 yields 433.
For example, the coefficient of
\[
J^{268}w^{39}s^{76}\eta^{37}i^8T^2\gamma hz
\]
is $-24\beta^{268}\ne0$. The displayed exponents sum to 433.
The delivered graph and independent leading-form calculation certify this
degree exactly, rather than merely bounding it from above.


\section{A ten-variable equation with a reused index carrier}
\label{sec:ten}

The following is a separate compression of the incoming ten-variable
product-square construction. Its main Pell base differs by one from the
base in Sections~\ref{sec:source}--\ref{sec:thirteen}; the definitions below
are complete and must be used together.

The variables are $J,F,\alpha,i,T,w,s,n,\eta,\zeta\in\Pos$.
Define
\begin{align*}
q&=\beta J+1,& h&=q-1,& C&=q-F-\alpha,\\
X&=wq,& Y&=q^3s,& D&=XY,\\
k&=\eta+\zeta,& c&=kY+\eta,& A&=D+Y+1,\\
\Delta&=A^2-1,& H&=4A-5,& V_0&=DY,\\
E&=1+V_0(V_0+1)k^2,\\
R&=[h(q-F)+\alpha](q^2-1)+(M_C+qM_F)J,\\
e&=k-R-1,& B&=(K_0+w)C-F,\\
U&=ic^2,& a&=\Delta U^2E,& F_0&=1+a,\\
G&=EF_0,& F_2&=1+4aF_0,& Q&=\Delta(F_2-1),\\
b&=cT-R,& V&=F_2b-c,& K_+&=X(3c-2X)+2.
\end{align*}
Choose $L=V^2$ for the 76-operation form, or $L=bV$ for the
77-operation form. Then put
\begin{align*}
M&=hDH,\\
Z&=hD\{H(n-LB)+K_+\}-H(DB+he),\\
R_G&=M^2G,\\
\boxed{P_{10}}&\boxed{=Q(Z^2-R_GV^2)+(R_G-Z^2)=0.}
\end{align*}
Every displayed symbol is polynomial; in particular, the proof-only rational
carrier below introduces neither a division gate nor a hidden variable.

\subsection{The new bounded-index recovery}

Write, only in this proof,
\[
z=Z/M,\quad \kappa=K_+/H,\quad \rho=e/D,\quad t=B/h,\quad W=1+hL.
\]
The exact identities are
\begin{equation}\label{eq:carrier}
P_{10}=(Q-1)Z^2-M^2G(QV^2-1),\qquad
z=n+\kappa-Wt-\rho.
\end{equation}
The inherited product-square extraction, with its signed adjacent-Pell
descent, first proves $G=g^2$ and $z=\pm gy\in\Z$.
It then recovers the main Pell index $R$ and the first pair
\[
c=Y_R(A),\qquad k=2Y_{P_0}(n_0),\qquad P_0=2XY^2+1,
\]
with $1\le n_0<R$ and $3q<R<q^4$. The sign arguments needed to justify
this order for the unshifted quotient are supplied below.

From the integrality of $z$, clearing only $hH$ in \eqref{eq:carrier}
shows $D\mid hHe$. Now
\[
q^4\mid D=XY=wq^4s,\qquad \gcd(q,hH)=1,
\]
because $h=q-1$ and $H\equiv-1\pmod q$.
Hence $q^4\mid e$, without assuming $\gcd(D,hH)=1$.
The elementary base-one Pell congruence gives
\[
k\equiv2n_0\pmod{4XY^2}.
\]
Consequently $q^4\mid 2n_0-R-1$, whereas
$|2n_0-R-1|<R<q^4$. Therefore $2n_0=R+1$ exactly.
Returning to the full congruence now proves $4XY^2\mid e$, hence
\begin{equation}\label{eq:recoverD}
D\mid e.
\end{equation}
The bounded index is authenticated before the full denominator is recovered.
Prime factors of $D$ shared with $hH$ never need to be cancelled.

Computationally $D$ is already useful in $A-1=D+Y$ and $V_0=DY$.
The three gates $XY,D+Y,DY$ replace four gates for
$X+1,Y(X+1),Y^2,XY^2$. This is the specific saved operation.

\subsection{Signed recovery makes the quotient unshifted}

The old shifted numerator used $b=c(T+\alpha)-R$. Omitting the addition
requires an argument that permits both signs of $b$ until the outer mask
is decoded. Positive inputs give $q\ge32$, $X\ge q$, $Y\ge q^3$,
$c\ge3$, and $F_2>c+1$. Thus $V\ne0$ for every integer $b$.
The identity $Q=4\Delta^2U^2G$ gives the same prime-valuation extraction
as before. If $z\ne0$, the positive rational number $|z|/g$ satisfies
\[
Q|V|^2-(Q-1)y^2=1.
\]
The adjacent-Pell descent makes $y$ integral: in lowest terms $y=r/m$,
$m^2\mid Q-1$; the transformation
\[
v'=(2Q-1)v-2(Q-1)y,\quad y'=(2Q-1)y-2Qv
\]
has a nonzero integral first coordinate with $|v'|<v$ for $v>1$ and
preserves the reduced denominator. Descent ends at $v=y=1$.
The equation itself rules out $z=0$ because $QV^2-1>0$.

The case $b=0$ would give
$0<y^2-c^2=(c^2-1)/(Q-1)<1$, impossible for integer $y$.
Elementary size bounds then give
\[
\operatorname{sign}V=\operatorname{sign}b,\quad
|V|>|b|,\quad |V|>3g,\quad y<\tfrac32|V|.
\]
The signed folding lemma supplies an odd main index $p$ and
$k=2Y_{P_0}(n_0)$. Excluding $n_0=1$ gives
$k>8XY^2$ and $c>k$. Thus $\kappa>0$.
For square transport $gy<V^2/2=L/2$ follows immediately.

For bilinear transport, the additional small-index exclusion is useful.
If $p=3$, write $u=X+2-2Y$. If $n_0\ge3$, Pell growth already gives
$kY>Y_A(3)=4A^2-1$, a contradiction. For $n_0=2$ one has
$k=8XY^2+4$, and exact expansion gives
\[
\eta=4Y^2(Xu+1)+8Y(X+1)+3-4Y.
\]
The inequality $0<\eta<k$ forces $u\in\{0,1\}$, so $2Y-X$ is one or
two. This contradicts $q\mid X,Y$ and $q\ge32$.
Hence $p\ge5$ and $c>4\Delta+3$.
The auxiliary odd index cannot be one or three: at three, the two signs
would give respectively
\[
F_2(b-4\Delta)=c-(4\Delta+3),\qquad
F_2(b+4\Delta)=c+(4\Delta+3),
\]
both impossible since their right sides lie strictly between zero and $F_2$.
Thus $|V|>8Q^2$, whence $|b|>4g$ and
$gy<3|bV|/8<L/2$.

In both variants $e=b-cT+k-1<b$, so $\rho<|b|/D$.
If $B\le-1$, equation \eqref{eq:carrier} gives
\[
n=\pm gy-\kappa+Wt+\rho<L/2-L+|b|/D<0,
\]
a contradiction. Therefore $B\ge0$, which forces $C>0$.
Only now apply the native mask inequality $3q<R<q^4$.
The lower bound on $c$ gives $c>2R$, hence $b=cT-R>0$.
The retained signed rank theorem then identifies $p=R$.
These steps justify the hypotheses used in the bounded-index argument.

\subsection{Completing the native decoding and the converse}

After \eqref{eq:recoverD}, the carrier gives $H\mid hK_+$.
The retained uniform ratio and Pell projection imply
\[
H\mid h(X-2^R)(2^RX+1),\qquad
|h(X-2^R)(2^RX+1)|<H.
\]
The whole product must vanish, so $X=2^R$; no nonunit factor is cancelled
through the composite modulus. Then $\kappa\in\Z$, the carrier gives
$Wt\in\Z$, and $W\equiv1\pmod h$ gives $h\mid B$.
Together with the literal $q^3\mid Y$, the inherited half-binomial rounding
and native reconstruction decode a finite refutation.

Conversely, canonical data from a refutation give $X=2^R$, $R\equiv3\pmod4$,
$q^3\mid Y$, and integral $t\ge0$. Choose a seed Pell ordinate divisible by
$c^2\tau$, double it, and use the uniform auxiliary reconstruction.
The recovered first index makes $\rho=e/(XY)>0$ integral, and the seed
can be chosen with $g>\kappa$. Thus
\[
n=gy+(1+hL)t+\rho-\kappa>0
\]
completes all ten positive coordinates. The full ordinary argument and its
inherited folding, ratio, and projection proofs are included in the
\file{compact_xy} component.

The exact degrees of the two supplied polynomials are 626 and 464.
For $L=V^2$, the unique leading summand is $QZ^2$; for $L=bV$ it is
$-QM^2GV^2$. Full leading homogeneous forms, rather than numerical samples,
certify the degrees in the supplied checkers.

\subsection{A lower-degree companion by delaying one factor of \texorpdfstring{$q$}{q}}

There is a further ten-positive tradeoff at 80 operations and exact degree
382. In the same definitions set $Y=q^2s$, retain $L=bV$, and replace the
carrier by
\begin{align*}
E_0&=q^4,& D_1&=q^5,& M&=hD_1H,\\
Z&=hD_1\{H(n-LB)+K_+\}-H\{D_1B+h(e-sE_0)\}.
\end{align*}
Use the same final polynomial $Q(Z^2-M^2GV^2)+(M^2G-Z^2)$.
As a proof identity, its carrier is now
\[
Z/M=n+\kappa+s/q-WB/h-e/D_1.
\]
The extra positive fraction $s/q$ is retained deliberately. It is represented
by the displayed ordinary polynomial numerator, not by an added operation
in the final equation.

The root and sign arguments still apply with $Y\ge q^2$. The initial lower
bound becomes $c>8q^5>2R$, and the exclusion of main index three still uses
only the literal $q\mid X,Y$. In the contradiction for negative $B$, the
new term contributes $-s/q<0$ and therefore strengthens the inequality.
Integrality gives $q^5\mid e-sq^4$, so $q^4\mid e$ before knowing $q\mid s$.
The same bounded-index argument proves $2n_0=R+1$. Then
$4XY^2\mid e$, and $q^5\mid XY^2=wq^5s^2$, so $D_1\mid e$.
The projection gives $X=2^R$ and integral $\kappa$ as before.
What remains is
\[
s/q-WB/h\in\Z.
\]
Clearing $qh$ and reducing separately modulo $q$ and $h$ proves
$q\mid s$ and $h\mid B$, since $\gcd(q,h)=1$ and $W\equiv1\pmod h$.
Thus the missing condition $q^3\mid Y$ is recovered exactly when the
native reconstruction needs it.

For completeness choose $s=Y/q^2$ from a canonical source witness; then
$s/q$ is integral. The fresh seed may be chosen with
$g>\kappa+s/q$, and
$n=gy+Wt+e/D_1-\kappa-s/q>0$ supplies the remaining coordinate.
This completes both directions. The literal graph, full proof, and exact
degree certificate are supplied as
\file{compact_xy/equation_80_382_partial.json} and
\file{compact_xy/PARTIAL_Q2_PROOF.md}.

\section{A direct quartic from a positive comparison center}
\label{sec:natural}

The direct construction is useful when avoiding the native K215 coefficient
tuple matters more than minimizing the number of variables. It starts from
the numerical proof calculus C222 and its retained equivalence to finite
ZF refutations and then finite ZFC refutations. No Jones parameter or K215
computation occurs in this branch.

The delivered principal polynomial is exactly
\begin{equation}\label{eq:quartic}
\boxed{P_{221}(x_0,\ldots,x_{220})=
\sum_{j=0}^{171}r_j(x_0,\ldots,x_{220})^2=0},
\qquad x_i\in\N.
\end{equation}
All 172 quadratic residuals are printed in Appendix~\ref{app:residuals}.
The machine-readable residual list and every collected coefficient are also
supplied. The displayed sum has a fixed upper limit, so it is a single
ordinary quartic, not a variable-size system or an unexpanded quantified
program.

\subsection{A positive polynomial center supplies two congruences}

An indexed Pell certificate has principal pair $(x,y)$ at base $A$, an
auxiliary pair $(u,v)$ at that base, and comparison pair $(s,r)$ at a base
$C$. The old source uses natural witnesses $d,e,f,g,h_0$ with
\begin{align*}
x^2-(A^2-1)y^2&=1,& u^2-(A^2-1)v^2&=1,\\
s^2-(C^2-1)r^2&=1,& v&=y^2(e+1),\\
C&=1+4y(d+1),& C&=A+uf,\\
s&=x+ug,&r&=H+4yh_0.
\end{align*}
The retained semantic cap gives $y\ge H\ge2$ before identifying the main
index. Replace the two definitions of $C$ and the auxiliary multiple by
\begin{equation}\label{eq:center}
\boxed{v=2y^2(e+1),\qquad C=A+(4y+1-A)u^2.}
\end{equation}
The quadratic graph uses the existing natural coordinate $C-1$ and a new
square wire for $u^2$; every intermediate quantified wire is counted.

\begin{lemma}[Positive center]\label{lem:center}
The modified indexed certificate is existence-equivalent to the original
one, in its retained semantic context. All restored old witnesses are natural.
\end{lemma}
\begin{proof}[Soundness]
Put $T=u^2$. The auxiliary norm gives
\[
T-1=4(A^2-1)y^4(e+1)^2,\qquad T>A.
\]
If $L=4y+1-A<0$, then $C=A+LT\le A-T<0$, contradicting the natural
coordinate $C-1$. Therefore $L\ge0$ and $C\ge A\ge2$.
Furthermore
\[
C-1=4yT-(A-1)(T-1)\equiv0\pmod{4y}.
\]
Restore
\[
d=(C-1)/(4y)-1\ge0,\quad f=Lu\ge0,\quad e_{\rm old}=2e+1\ge0.
\]
All old rows hold. This argument never presupposes $y=Y_H(A)$.
\end{proof}
\begin{proof}[Completeness]
For canonical principal witnesses $y=Y_H(A)$ with $H\ge2$, one has
$y\ge2A$. Choose a positive auxiliary ordinate $v$ divisible by $2y^2$
using Lemma~\ref{lem:period}. Then $e=v/(2y^2)-1\ge0$.
The center in \eqref{eq:center} has $C>A$, $C\equiv A\pmod u$,
and $C\equiv1\pmod{4y}$. Take $s=X_H(C)$ and $r=Y_H(C)$.
These are integer polynomials in their bases, so
\[
s\equiv X_H(A)=x\pmod u,\qquad r\equiv Y_H(1)=H\pmod{4y}.
\]
Monotonicity gives $s>x$ and $r\ge H$, so both remaining quotients
are natural. The principal pair, seed pairs, and semantic-cap witness
are unchanged.
\end{proof}

The two arbitrary congruence parameters have been replaced by one fixed
polynomial expression. The natural domain supplies the sign needed for the
inverse map. The new square wire has a cost; after exact elimination, this
center saves one net coordinate in the selected implementation.

\subsection{Deleting a bound before using its intended interpretation}

The retained source has $t\ge12$, $K=x\ge F+1$, and the six exact-power
certificates including
\[
V_0=(F^2+1)^t,\qquad N=B^K,\qquad q=(NS)^K.
\]
The powers are established from the unchanged frame and cap before any
factorial or binomial interpretation. The ceiling rows imply
\[
V_0+s_f=BF.
\]
As $V_0>F^2$, we obtain $B>F\ge1$ and $K\ge2$. Hence
\[
B<B^2\le N\le q.
\]
The private bound row $q=B+o+1$ can therefore be removed and restored later
by $o=q-B-1\ge0$. This deletion does not use $F=t!$ or the intended binomial
meaning of $B$. Avoiding that circularity is essential.

\subsection{Pooling eleven remainder bounds into four}

For a group with common denominator $D$, replace the separate rows
$D=r_i+s_i+1$ by
\[
D=\sum_i r_i+s+1,\qquad s\in\N.
\]
Soundness is immediate: restore each old slack as
$s_i=s+\sum_{j\ne i}r_j$. Completeness is stronger than the old separate
bounds and requires a genuine aggregate estimate.

The four denominators in the actual graph are
$q^t,q^{t+1},q^{y_\beta+1},q^x$. The first has two requests, the second
five, and the others two each. In canonical completeness witnesses, the
self-sized radix has $q=(NS)^K$, $K=x$, all ordinary upper indices are
less than $S$, and the derived label satisfies $J\le S^{2^{22}}$ with
$2^{22}(t+1)<F<K$. Every ordinary or label low-digit coefficient is below
$S^K$. For the factorial upper index $F^2+t$, canonical $F=t!$ gives
$F^2+t\ge2t$ and
\[
\binom{F^2+t}{j}\le\binom{F^2+t}{t}
=\prod_{r=1}^t\frac{F^2+r}{r}\le(F^2+1)^t=V_0\le N
\quad(0\le j\le t).
\]
Thus each group's coefficient sum is less than
\[
5S^K+N<N^KS^K=q.
\]
For the final inequality use $N\ge4$, $K\ge2$, $S\ge2$, and
$N^K-5\ge N$. Each low remainder is a sum of these coefficients times
the corresponding powers of $q$. The aggregate remainder is therefore
at most $q^k-1$ for denominator $q^k$, supplying the new natural slack.
Factorial semantics are used here only in the construction of already
canonical completeness witnesses.

\subsection{From the raw graph to the literal quartic}

The modified raw graph has 324 natural coordinates and 275 quadratic
residuals. Exact monic substitutions remove 103 coordinates for the principal
form, or 102 for the sparse companion. Every substitution has an explicit
polynomial inverse; nonnegativity is certified by rational sums of squares
plus nonnegative monomials. These certificates are replayed independently.
The final equations have respectively 221/172 and 222/173 variables/residuals.

Soundness first restores the pooled slacks, then the indexed Pell certificate,
then every exact power, and only then the deleted digit-output bound.
The entire original arithmetic source is reconstructed at that point.
Its retained C222/refutation interfaces finish the implication.
Completeness starts with a finite refutation, chooses the canonical semantic
and CRT data, supplies the aggregate slacks, chooses the fresh positive center,
and projects by the certified monic substitutions. Both directions preserve
the explicitly stated natural domain.

\section{Verification, limits, and how to use the archive}
\label{sec:verification}

The archive is intended for review and incorporation as a new ProveIt intake
package. It does not change upstream branches or replace earlier constructions.
The principal components are:
\begin{itemize}
\item \file{products/}: literal thirteen- and fourteen-positive circuits,
ordinary proofs, source binding, and exact circuit checks.
\item \file{compact_xy/}: the ten-positive carrier circuits, their full signed
proof, literal coefficients, and independent coefficient reconstruction.
\item \file{natural/}: the two direct quartics, every expanded coefficient,
the raw graph, all restoration certificates, full new proofs, the retained
arithmetic producer closure, selected logical interface sources, and pinned
calibration references.
\item \file{report/}: this editable report, complete fixed-constant recipe,
and the printed 172-residual appendix.
\end{itemize}

The Python reproduction programs use the standard library and do not need
an external checkout or network connection. The top-level verifier runs the
component checks; its optional rebuild mode regenerates the direct outputs
in temporary directories and compares the specified mathematical artifacts.
The supplied receipts record actual invocations. The TeX report compiles with
pdfLaTeX using ordinary TeX Live packages.

The finite checks include literal graph topology, liveness of every variable,
exact local polynomial identities, complete leading homogeneous forms,
fixed coefficient bytes, all 275 raw direct rows, exact monic-restoration
certificates, all final quartic coefficients, and parsing of every printed
direct residual back into its coefficient data. Deliberate corruptions and
small exact Pell examples exercise the implementation. Such examples are
diagnostics; the universal proofs in the preceding sections do the logical work.

The new ordinary reductions were independently reviewed within the assisted
research workflow. That is not outside human peer review. The inherited
native compiler and C222/ZF/ZFC calibrations are explicitly retained rather
than silently replaced by finite testing. No new end-to-end Lean/Rocq build,
encoded object-language ZFC proof, global operation lower bound, or decision
of consistency is claimed. A future formalization can profitably start with
Lemma~\ref{lem:product}, the two product-certificate restoration maps, and the
positive-center lemma: these are small, separable arithmetic targets.

\appendix
\clearpage
\section{Complete finite recipe for the compact coefficient tuple}
\label{app:constants}

The following retained recipe fixes all four positive integer ports used by
the compact constructions. It is the entire indexed943/radix67 tuple.
In the report formulas $K_0=K$. All indices in this appendix are bounded
constant indices; none is an additional Diophantine unknown. The floors,
moduli, finite sums, and rational intermediate below define fixed integers
only. The final equations themselves use solely polynomial operations.
The component reconstruction program checks every canonical coefficient
byte against this recipe. The largest port has 592,884 bits, which should
not be confused with the coefficient height of the fully expanded polynomial.

\begingroup\small
\subsection*{The four fixed integers}
Every index, set, sum, product and floor below defines a constant, not
another unknown. Put
\[
\rho=2^{67},\quad B=\rho^{8807},\quad\beta=B-1,\quad
G_m(x)=\sum_{j=0}^{m-1}x^j.
\]
Define
\[
\begin{split}
\mathcal E={}&\rg0{942}\cup\rg{983}{992}\cup\rg{1003}{1011}\\
&\cup\rg{1013}{1031}\cup\{1063,1064,1095,2195,3296\}.
\end{split}
\]
Then
\begin{align*}
M_C={}&B-1-\sum_{e\in\mathcal E\setminus\{1\}}\rho^e-2\rho^{1063},\\
M_F={}&B-1+4+(2^{51}-2)(\rho^{4393}+\rho^{5495})\\
&+\rho^{1013}G_{10}(\rho)+\rho^{1065}G_9(\rho)
  +\rho^{2195}+\rho^{3296}\\
&+(2^{65}-1)\rho^{8792}G_{13}(\rho)+(2^{19}-1)\rho^{8805},\\
K={}&\rho^{30}+\rho^{62}+\rho^{1100}+\rho^{2201}+\rho^{8849}\\
&+(\rho^{4393}+\rho^{5495})\Phi(1/\rho),
\end{align*}
where
\[
\begin{split}
\Phi(x)={}&(1+2^{51})G_{943}(x)
 -2^{10}\sum_{i=0}^{942}v(t_i)x^i\\
&+\sum_{(b,e,m)\in\mathcal C}2^b x^e G_m(2x),\\
\mathcal C={}&\{(10,1023,8),(18,1013,9),(27,1003,8),\\
&(35,983,9),(44,1031,1),(45,1022,1),(46,1011,1),\\
&(47,992,1),(48,1095,1),(49,2195,1),(50,3296,1)\}.
\end{split}
\]
The rational intermediate in K cancels exactly; all four results are positive integers.

\subsection*{The 943 indexed fixed tiles}
Use the actions $u_j,v_j,d_j$ decoded below. Let
\[
\begin{gathered}
h(q,b)=2(q+1)+b,\quad
\mathcal R=\{u_j:d_j=1,\ 0\le j<430\},\\
\mathcal L=\{u_j:d_j=0,\ 0\le j<430\},\\
r_q=6+|\mathcal R\cap\rg0{q-1}|,\quad
l_q=205+|\mathcal L\cap\rg0{q-1}|.
\end{gathered}
\]
An interval ending at $-1$ is empty. The table defines $t_i$ for
$0\le i\le942$; all index ranges are fixed:
\[
\setlength{\arraycolsep}{1.4pt}\begin{array}{c|l}
\text{index }i&\text{tile }t_i=(w,n,e,s)\\\hline
0,1&(2,435,3,1),\ (5,0,0,0)\\
j+2\ (j=2q+b<430)&((1-d_j)l_{u_j},h(q,b),d_jr_{u_j},v_j)\\
432+b\ (b=0,1)&(0,h(215,b),0,h(215,b))\\
434+b\ (b=0,1)&(0,b,0,b)\\
424+2r_q+b\ (q\in\mathcal R,b=0,1)&(r_q,b,0,h(q,b))\\
424+2l_q+b\ (q\in\mathcal L,b=0,1)&(0,b,l_q,h(q,b))\\
924+j\ (0\le j<4)&(0,j+430\fl{j/2},0,438)\\
928&(2,434,2,0)\\
929&(3,436,4,4)\\
930&(4,436,4,0)\\
497+t\ (t=434,435,436)&(1,438,1,t)\\
\substack{934+3a+b-6\fl{a/4}\\a\in\{0,1,4\},\ b=0,1,2}&(a,437,b,437)
\end{array}
\]
For $t=(w,n,e,s)$ put
\[
v(t)=w+2^8n+2^{17}e+2^{25}s+127\,2^{34}.
\]

\subsection*{A closed decoder for the 430 actions}
Put $D_b(A,k)=\fl{A/b^k}\bmod b$, where
$a\bmod b=a-b\fl{a/b}$. For $0\le q<215$, $b=0,1$, $j=2q+b$:
\[
\begin{gathered}
N=\mathcal A\bmod2^{636},\quad
T_0=\fl{\mathcal A/2^{636}}\bmod2^{1675},\\
E_0=\fl{\mathcal A/2^{2311}},\\
\nu_j=D_2(N,j),\quad \sigma_j=\sum_{k=0}^{j-1}\nu_k,\\
m_q=D_2(N,430+q),\quad \tau_q=\sum_{k=0}^{q-1}m_k,\\
F_q=14(1-m_q)+m_qD_{16}(E_0,\tau_q),\\
u_j=\nu_j(1+\sigma_j)+(1-\nu_j)D_{216}(T_0,j-\sigma_j),\\
v_j=D_2(F_q,2b),\quad d_j=D_2(F_q,2b+1).
\end{gathered}
\]
Empty sums are zero. Here $(v_0,d_0,u_0)=(1,1,1)$; the constructor tile
uses this actual first target, including its state renaming.

\subsection*{The integer $\mathcal A$}
The alphabet, in order of digit values $0,\ldots,61$, is
{\fontsize{7.7}{8.2}\selectfont\ttfamily 0123456789ABCDEFGHIJKLMNOPQRSTUVWXYZabcdefghijklmnopqrstuvwxyz}.\par
Concatenate the following 436 characters as digits $e_0,\ldots,e_{435}$;
put $\mathcal A=\sum_{j=0}^{435}e_j62^{435-j}$.\par
\begingroup\fontsize{7.7}{8.2}\selectfont\ttfamily
BxEnhvy6NipzkWmROo0GsA3uK3twWLxscLq20UTumy5Cdu461LwD6g7RDNzF8Y\par
iYDFTpB8ehAypd7d6OsOL4APamnJeNXQypUezCmW9n4oXP6XfgmiWuAWZ8ZJ7U\par
OP7gu4kVdlGZqZHULQMar1dXNeJW4i7oNVI8usp8BlovjjZbwPCovZI328fRUq\par
4q92R81s5GjWQaQfGhkq9UYMK3rRMaBmsOIBkJqc2cj2VRIZ6mjzkBIOkW9SjZ\par
FzeA6J5u4j3sBL3SdgwXzM9mHvQivu8EJ1TBTpRhHDqZIbmYgwQ8VqTAFotl6F\par
We1YC554GKrVHsA8oUGKSA6F6k5KREBSD0Tgd7bFALY84LLbpw574BQyvdoHrl\par
0Oe35uwDf15DpT6VJWZBjSbHPpPAINYTte0eUxnTFTKBlGj8Am6sb1gLldh8WZ\par
8P\par
\endgroup

\endgroup

\clearpage
\section{All residuals of the 221-variable quartic}
\label{app:residuals}

For $x_0,\ldots,x_{220}\in\N$, the equation is exactly
$\sum_{j=0}^{171}r_j^2=0$, using the following complete list.
There are no additional definitions or quantified registers in this appendix.
The transcription is parsed back to the literal sparse residual data by
\file{natural/render.py --check}.
\begingroup\small
\setlength{\abovedisplayskip}{4pt}\setlength{\belowdisplayskip}{4pt}
% Generated from the complete literal equation.json; no arithmetic macros.
\begin{align*}
r_{0}&=
-x_{2}^{2}+x_{3}
\end{align*}
\begin{align*}
r_{1}&=
-x_{2}-x_{2}x_{5}-x_{2}x_{6}+x_{4}+x_{6}
\end{align*}
\begin{align*}
r_{2}&=
-13-x_{0}-x_{0}x_{2}-x_{0}x_{8}-13x_{2}-13x_{8}+x_{10}
\end{align*}
\begin{align*}
r_{3}&=
-1+x_{1}-x_{2}-x_{2}x_{7}-2x_{7}-x_{7}x_{8}
\end{align*}
\begin{align*}
r_{4}&=
-1+x_{1}-x_{9}-x_{9}x_{10}
\end{align*}
\begin{align*}
r_{5}&=
2+x_{3}-x_{5}x_{21}-x_{6}x_{21}+x_{19}-x_{21}
\end{align*}
\begin{align*}
r_{6}&=
-6-5x_{11}-x_{11}^{2}-2x_{11}x_{12}-5x_{12}-x_{12}^{2}+x_{22}
\end{align*}
\begin{align*}
r_{7}&=
-72-68x_{12}-16x_{12}^{2}-16x_{12}x_{22}-34x_{22}-4x_{22}^{2}+x_{23}
\end{align*}
\begin{align*}
r_{8}&=
-30-11x_{11}-x_{11}^{2}-4x_{11}x_{23}-22x_{23}-4x_{23}^{2}+x_{24}
\end{align*}
\begin{align*}
r_{9}&=
-380-312x_{23}-64x_{23}^{2}-32x_{23}x_{24}-78x_{24}-4x_{24}^{2}+x_{25}
\end{align*}
\begin{align*}
r_{10}&=
-6-5x_{12}-x_{12}^{2}-2x_{12}x_{13}-5x_{13}-x_{13}^{2}+x_{26}
\end{align*}
\begin{align*}
r_{11}&=
-6-5x_{11}-x_{11}^{2}-2x_{11}x_{13}-5x_{13}-x_{13}^{2}+x_{27}
\end{align*}
\begin{align*}
r_{12}&=
-414-22x_{11}-2x_{11}^{2}-16x_{11}x_{13}-8x_{11}x_{26}-200x_{12}-32x_{12}^{2}-64x_{12}x_{13}-32x_{12}x_{22} \\
&\quad -32x_{12}x_{27}-320x_{13}-64x_{13}^{2}-32x_{13}x_{22}-32x_{13}x_{26}-32x_{13}x_{27}-100x_{22}-8x_{22}^{2} \\
&\quad -16x_{22}x_{27}-52x_{26}-8x_{26}^{2}-108x_{27}-8x_{27}^{2}+x_{28}
\end{align*}
\begin{align*}
r_{13}&=
-1344-800x_{12}-128x_{12}^{2}-256x_{12}x_{13}-128x_{12}x_{22}-128x_{12}x_{27}-864x_{13}-128x_{13}^{2} \\
&\quad -128x_{13}x_{22}-128x_{13}x_{27}-400x_{22}-32x_{22}^{2}-64x_{22}x_{27}-432x_{27}-32x_{27}^{2}-2x_{28} \\
&\quad -2x_{28}^{2}+x_{29}
\end{align*}
\begin{align*}
r_{14}&=
-247-168x_{11}-32x_{11}^{2}-64x_{11}x_{12}-32x_{11}x_{26}-32x_{11}x_{27}-188x_{12}-32x_{12}^{2}-32x_{12}x_{26} \\
&\quad -32x_{12}x_{27}-2x_{22}-92x_{26}-8x_{26}^{2}-16x_{26}x_{27}-84x_{27}-8x_{27}^{2}+x_{30}
\end{align*}
\begin{align*}
r_{15}&=
-960-672x_{11}-128x_{11}^{2}-256x_{11}x_{12}-128x_{11}x_{26}-128x_{11}x_{27}-736x_{12}-128x_{12}^{2} \\
&\quad -128x_{12}x_{26}-128x_{12}x_{27}-368x_{26}-32x_{26}^{2}-64x_{26}x_{27}-336x_{27}-32x_{27}^{2}-2x_{30} \\
&\quad -2x_{30}^{2}+x_{31}
\end{align*}
\begin{align*}
r_{16}&=
-276-30x_{11}-2x_{11}^{2}-16x_{11}x_{13}-8x_{11}x_{26}-30x_{12}-2x_{12}^{2}-16x_{12}x_{13}-8x_{12}x_{27} \\
&\quad -272x_{13}-64x_{13}^{2}-32x_{13}x_{26}-32x_{13}x_{27}-68x_{26}-8x_{26}^{2}-68x_{27}-8x_{27}^{2}+x_{32}
\end{align*}
\begin{align*}
r_{17}&=
-552-120x_{12}-8x_{12}^{2}-64x_{12}x_{13}-32x_{12}x_{27}-544x_{13}-128x_{13}^{2}-128x_{13}x_{27}-272x_{27} \\
&\quad -32x_{27}^{2}-2x_{32}-2x_{32}^{2}+x_{33}
\end{align*}
\begin{align*}
r_{18}&=
-247-4x_{11}-352x_{12}-128x_{12}^{2}-64x_{12}x_{22}-64x_{12}x_{26}-92x_{22}-8x_{22}^{2}-16x_{22}x_{26}-84x_{26} \\
&\quad -8x_{26}^{2}-2x_{27}+x_{34}
\end{align*}
\begin{align*}
r_{19}&=
-960-1408x_{12}-512x_{12}^{2}-256x_{12}x_{22}-256x_{12}x_{26}-368x_{22}-32x_{22}^{2}-64x_{22}x_{26}-336x_{26} \\
&\quad -32x_{26}^{2}-2x_{34}-2x_{34}^{2}+x_{35}
\end{align*}
\begin{align*}
r_{20}&=
-56-60x_{12}-16x_{12}^{2}-16x_{12}x_{22}-30x_{22}-4x_{22}^{2}+x_{36}
\end{align*}
\begin{align*}
r_{21}&=
-30-11x_{11}-x_{11}^{2}-4x_{11}x_{36}-22x_{36}-4x_{36}^{2}+x_{37}
\end{align*}
\begin{align*}
r_{22}&=
-380-312x_{36}-64x_{36}^{2}-32x_{36}x_{37}-78x_{37}-4x_{37}^{2}+x_{38}
\end{align*}
\begin{align*}
r_{23}&=
-30-11x_{12}-x_{12}^{2}-4x_{12}x_{38}-22x_{38}-4x_{38}^{2}+x_{39}
\end{align*}
\begin{align*}
r_{24}&=
-349-420x_{12}-128x_{12}^{2}-128x_{12}x_{22}-208x_{22}-32x_{22}^{2}-2x_{26}+x_{40}
\end{align*}
\begin{align*}
r_{25}&=
-1368-1664x_{12}-512x_{12}^{2}-512x_{12}x_{22}-832x_{22}-128x_{22}^{2}-2x_{40}-2x_{40}^{2}+x_{41}
\end{align*}
\begin{align*}
r_{26}&=
-56-60x_{11}-16x_{11}^{2}-16x_{11}x_{27}-8x_{11}x_{41}-30x_{27}-4x_{27}^{2}-4x_{27}x_{41}-15x_{41}-x_{41}^{2} \\
&\quad +x_{42}
\end{align*}
\begin{align*}
r_{27}&=
-6-5x_{13}-x_{13}^{2}-2x_{13}x_{14}-5x_{14}-x_{14}^{2}+x_{43}
\end{align*}
\begin{align*}
r_{28}&=
-224-26x_{11}-2x_{11}^{2}-16x_{11}x_{14}-8x_{11}x_{43}-288x_{13}-128x_{13}^{2}-64x_{13}x_{26}-64x_{13}x_{27} \\
&\quad -88x_{14}-32x_{14}^{2}-32x_{14}x_{43}-76x_{26}-8x_{26}^{2}-16x_{26}x_{27}-68x_{27}-8x_{27}^{2}-44x_{43} \\
&\quad -8x_{43}^{2}+x_{44}
\end{align*}
\begin{align*}
r_{29}&=
-640-1152x_{13}-512x_{13}^{2}-256x_{13}x_{26}-256x_{13}x_{27}-304x_{26}-32x_{26}^{2}-64x_{26}x_{27}-272x_{27} \\
&\quad -32x_{27}^{2}-2x_{44}-2x_{44}^{2}+x_{45}
\end{align*}
\begin{align*}
r_{30}&=
-241-168x_{11}-32x_{11}^{2}-64x_{11}x_{12}-32x_{11}x_{26}-32x_{11}x_{27}-184x_{12}-32x_{12}^{2}-32x_{12}x_{26} \\
&\quad -32x_{12}x_{27}-x_{13}-92x_{26}-8x_{26}^{2}-16x_{26}x_{27}-84x_{27}-8x_{27}^{2}+x_{46}
\end{align*}
\begin{align*}
r_{31}&=
-42-52x_{13}-16x_{13}^{2}-16x_{13}x_{26}-26x_{26}-4x_{26}^{2}+x_{47}
\end{align*}
\begin{align*}
r_{32}&=
-30-11x_{13}-x_{13}^{2}-4x_{13}x_{47}-22x_{47}-4x_{47}^{2}+x_{48}
\end{align*}
\begin{align*}
r_{33}&=
-306-35x_{12}-x_{12}^{2}-16x_{12}x_{47}-4x_{12}x_{48}-280x_{47}-64x_{47}^{2}-32x_{47}x_{48}-70x_{48}-4x_{48}^{2} \\
&\quad +x_{49}
\end{align*}
\begin{align*}
r_{34}&=
-4556-4320x_{47}-1024x_{47}^{2}-512x_{47}x_{48}-128x_{47}x_{49}-1080x_{48}-64x_{48}^{2}-32x_{48}x_{49}-270x_{49} \\
&\quad -4x_{49}^{2}+x_{50}
\end{align*}
\begin{align*}
r_{35}&=
-30-11x_{13}-x_{13}^{2}-4x_{13}x_{50}-22x_{50}-4x_{50}^{2}+x_{51}
\end{align*}
\begin{align*}
r_{36}&=
-306-35x_{12}-x_{12}^{2}-16x_{12}x_{50}-4x_{12}x_{51}-280x_{50}-64x_{50}^{2}-32x_{50}x_{51}-70x_{51}-4x_{51}^{2} \\
&\quad +x_{52}
\end{align*}
\begin{align*}
r_{37}&=
-1028-672x_{11}-128x_{11}^{2}-256x_{11}x_{12}-128x_{11}x_{26}-128x_{11}x_{27}-736x_{12}-128x_{12}^{2} \\
&\quad -128x_{12}x_{26}-128x_{12}x_{27}-368x_{26}-32x_{26}^{2}-64x_{26}x_{27}-336x_{27}-32x_{27}^{2}-2x_{46} \\
&\quad -2x_{46}^{2}-32x_{50}-8x_{51}-2x_{52}+x_{53}
\end{align*}
\begin{align*}
r_{38}&=
-264-128x_{50}-32x_{51}-8x_{52}-2x_{53}-2x_{53}^{2}+x_{54}
\end{align*}
\begin{align*}
r_{39}&=
-2-3x_{11}-x_{11}^{2}-2x_{11}x_{54}-3x_{54}-x_{54}^{2}+x_{55}
\end{align*}
\begin{align*}
r_{40}&=
-12-28x_{54}-16x_{54}^{2}-16x_{54}x_{55}-14x_{55}-4x_{55}^{2}+x_{56}
\end{align*}
\begin{align*}
r_{41}&=
-30-11x_{13}-x_{13}^{2}-4x_{13}x_{56}-22x_{56}-4x_{56}^{2}+x_{57}
\end{align*}
\begin{align*}
r_{42}&=
-306-35x_{11}-x_{11}^{2}-16x_{11}x_{56}-4x_{11}x_{57}-280x_{56}-64x_{56}^{2}-32x_{56}x_{57}-70x_{57}-4x_{57}^{2} \\
&\quad +x_{58}
\end{align*}
\begin{align*}
r_{43}&=
-4556-4320x_{56}-1024x_{56}^{2}-512x_{56}x_{57}-128x_{56}x_{58}-1080x_{57}-64x_{57}^{2}-32x_{57}x_{58}-270x_{58} \\
&\quad -4x_{58}^{2}+x_{59}
\end{align*}
\begin{align*}
r_{44}&=
-90-76x_{13}-16x_{13}^{2}-16x_{13}x_{26}-16x_{13}x_{59}-38x_{26}-4x_{26}^{2}-8x_{26}x_{59}-38x_{59}-4x_{59}^{2} \\
&\quad +x_{60}
\end{align*}
\begin{align*}
r_{45}&=
-4556-4320x_{50}-1024x_{50}^{2}-512x_{50}x_{51}-128x_{50}x_{52}-1080x_{51}-64x_{51}^{2}-32x_{51}x_{52}-270x_{52} \\
&\quad -4x_{52}^{2}+x_{61}
\end{align*}
\begin{align*}
r_{46}&=
-682-70x_{12}-2x_{12}^{2}-32x_{12}x_{59}-8x_{12}x_{60}-592x_{59}-128x_{59}^{2}-64x_{59}x_{60}-148x_{60} \\
&\quad -8x_{60}^{2}-2x_{61}+x_{62}
\end{align*}
\begin{align*}
r_{47}&=
-x_{62}-x_{62}^{2}+x_{63}
\end{align*}
\begin{align*}
r_{48}&=
-42-52x_{12}-16x_{12}^{2}-8x_{12}x_{13}-16x_{12}x_{22}-13x_{13}-x_{13}^{2}-4x_{13}x_{22}-26x_{22}-4x_{22}^{2} \\
&\quad +x_{64}
\end{align*}
\begin{align*}
r_{49}&=
-2278-2264x_{12}-800x_{12}^{2}-800x_{12}x_{22}-160x_{12}x_{64}-70x_{13}-2x_{13}^{2}-32x_{13}x_{61}-8x_{13}x_{63} \\
&\quad -1132x_{22}-200x_{22}^{2}-80x_{22}x_{64}-592x_{61}-128x_{61}^{2}-64x_{61}x_{63}-148x_{63}-8x_{63}^{2} \\
&\quad -228x_{64}-8x_{64}^{2}+x_{65}
\end{align*}
\begin{align*}
r_{50}&=
-2712-280x_{13}-8x_{13}^{2}-128x_{13}x_{61}-32x_{13}x_{63}-2368x_{61}-512x_{61}^{2}-256x_{61}x_{63}-592x_{63} \\
&\quad -32x_{63}^{2}-2x_{65}-2x_{65}^{2}+x_{66}
\end{align*}
\begin{align*}
r_{51}&=
-2-3x_{66}-x_{66}^{2}+x_{67}
\end{align*}
\begin{align*}
r_{52}&=
-x_{2}-x_{2}x_{16}+x_{68}
\end{align*}
\begin{align*}
r_{53}&=
-x_{15}^{2}+2x_{15}x_{16}-x_{16}^{2}+x_{69}
\end{align*}
\begin{align*}
r_{54}&=
-2x_{25}x_{68}-4x_{68}+x_{70}
\end{align*}
\begin{align*}
r_{55}&=
-x_{15}^{2}+2x_{15}x_{16}+2x_{15}x_{70}-x_{16}^{2}-2x_{16}x_{70}-x_{70}^{2}+x_{71}
\end{align*}
\begin{align*}
r_{56}&=
-x_{29}x_{68}+x_{72}
\end{align*}
\begin{align*}
r_{57}&=
-x_{15}^{2}+2x_{15}x_{16}+2x_{15}x_{72}-x_{16}^{2}-2x_{16}x_{72}-x_{72}^{2}+x_{73}
\end{align*}
\begin{align*}
r_{58}&=
-x_{31}x_{68}+x_{74}
\end{align*}
\begin{align*}
r_{59}&=
-x_{15}^{2}+2x_{15}x_{16}+2x_{15}x_{74}-x_{16}^{2}-2x_{16}x_{74}-x_{74}^{2}+x_{75}
\end{align*}
\begin{align*}
r_{60}&=
-x_{33}x_{68}+x_{76}
\end{align*}
\begin{align*}
r_{61}&=
-x_{15}^{2}+2x_{15}x_{16}+2x_{15}x_{76}-x_{16}^{2}-2x_{16}x_{76}-x_{76}^{2}+x_{77}
\end{align*}
\begin{align*}
r_{62}&=
-x_{35}x_{68}+x_{78}
\end{align*}
\begin{align*}
r_{63}&=
-x_{15}^{2}+2x_{15}x_{16}+2x_{15}x_{78}-x_{16}^{2}-2x_{16}x_{78}-x_{78}^{2}+x_{79}
\end{align*}
\begin{align*}
r_{64}&=
-4x_{12}x_{68}-2x_{39}x_{68}-4x_{68}+x_{80}
\end{align*}
\begin{align*}
r_{65}&=
-x_{15}^{2}+2x_{15}x_{16}+2x_{15}x_{80}-x_{16}^{2}-2x_{16}x_{80}-x_{80}^{2}+x_{81}
\end{align*}
\begin{align*}
r_{66}&=
-4x_{41}x_{68}-2x_{42}x_{68}+x_{82}
\end{align*}
\begin{align*}
r_{67}&=
-x_{15}^{2}+2x_{15}x_{16}+2x_{15}x_{82}-x_{16}^{2}-2x_{16}x_{82}-x_{82}^{2}+x_{83}
\end{align*}
\begin{align*}
r_{68}&=
-x_{45}x_{68}+x_{84}
\end{align*}
\begin{align*}
r_{69}&=
-x_{15}^{2}+2x_{15}x_{16}+2x_{15}x_{84}-x_{16}^{2}-2x_{16}x_{84}-x_{84}^{2}+x_{85}
\end{align*}
\begin{align*}
r_{70}&=
-x_{11}^{2}+2x_{11}x_{13}-x_{13}^{2}+x_{86}
\end{align*}
\begin{align*}
r_{71}&=
-x_{12}^{2}+2x_{12}x_{13}-x_{13}^{2}+x_{87}
\end{align*}
\begin{align*}
r_{72}&=
-x_{86}x_{87}+x_{88}
\end{align*}
\begin{align*}
r_{73}&=
-4-4x_{14}-x_{14}^{2}+x_{89}
\end{align*}
\begin{align*}
r_{74}&=
-x_{88}^{2}+2x_{88}x_{89}-x_{89}^{2}+x_{90}
\end{align*}
\begin{align*}
r_{75}&=
-2x_{67}x_{68}-4x_{68}+x_{91}
\end{align*}
\begin{align*}
r_{76}&=
-x_{15}^{2}+2x_{15}x_{16}+2x_{15}x_{91}-x_{16}^{2}-2x_{16}x_{91}-x_{91}^{2}+x_{92}
\end{align*}
\begin{align*}
r_{77}&=
-x_{2}-x_{2}x_{12}+x_{93}
\end{align*}
\begin{align*}
r_{78}&=
-1-x_{16}-x_{16}x_{93}-x_{93}+x_{94}
\end{align*}
\begin{align*}
r_{79}&=
-x_{2}-x_{2}x_{11}+x_{95}
\end{align*}
\begin{align*}
r_{80}&=
-4x_{2}-4x_{2}x_{12}-2x_{2}x_{22}+x_{96}
\end{align*}
\begin{align*}
r_{81}&=
-1-x_{15}-x_{15}x_{95}-x_{95}+x_{97}
\end{align*}
\begin{align*}
r_{82}&=
-x_{96}x_{97}-x_{97}+x_{98}
\end{align*}
\begin{align*}
r_{83}&=
-x_{94}^{2}+2x_{94}x_{98}-x_{98}^{2}+x_{99}
\end{align*}
\begin{align*}
r_{84}&=
-6x_{2}-4x_{2}x_{12}-2x_{2}x_{22}+x_{100}
\end{align*}
\begin{align*}
r_{85}&=
-1-x_{16}-x_{16}x_{100}-x_{100}+x_{101}
\end{align*}
\begin{align*}
r_{86}&=
-1-x_{15}-x_{15}x_{93}-x_{93}+x_{102}
\end{align*}
\begin{align*}
r_{87}&=
-x_{101}^{2}+2x_{101}x_{102}-x_{102}^{2}+x_{103}
\end{align*}
\begin{align*}
r_{88}&=
-x_{69}x_{71}+x_{104}
\end{align*}
\begin{align*}
r_{89}&=
-x_{73}x_{104}+x_{105}
\end{align*}
\begin{align*}
r_{90}&=
-x_{75}x_{105}+x_{106}
\end{align*}
\begin{align*}
r_{91}&=
-x_{77}x_{106}+x_{107}
\end{align*}
\begin{align*}
r_{92}&=
-x_{79}x_{107}+x_{108}
\end{align*}
\begin{align*}
r_{93}&=
-x_{81}x_{108}+x_{109}
\end{align*}
\begin{align*}
r_{94}&=
-x_{83}x_{109}+x_{110}
\end{align*}
\begin{align*}
r_{95}&=
-x_{85}x_{110}+x_{111}
\end{align*}
\begin{align*}
r_{96}&=
-x_{90}x_{111}-x_{92}x_{111}+x_{112}
\end{align*}
\begin{align*}
r_{97}&=
-x_{99}x_{112}+x_{113}
\end{align*}
\begin{align*}
r_{98}&=
-2-2x_{2}-x_{2}x_{19}-2x_{8}-x_{8}x_{19}-x_{19}+x_{114}
\end{align*}
\begin{align*}
r_{99}&=
-1-x_{17}-x_{17}x_{114}-x_{114}+x_{115}
\end{align*}
\begin{align*}
r_{100}&=
-2-x_{2}-x_{2}x_{18}-x_{8}-x_{8}x_{18}-2x_{18}-x_{18}x_{114}-x_{114}+x_{116}
\end{align*}
\begin{align*}
r_{101}&=
2-4x_{1}+2x_{1}^{2}-2x_{1}x_{15}-2x_{1}x_{16}-2x_{1}x_{115}-2x_{1}x_{116}+2x_{15}+x_{15}^{2}+2x_{15}x_{115} \\
&\quad +2x_{16}+x_{16}^{2}+2x_{16}x_{116}-x_{20}x_{117}+x_{103}x_{113}+2x_{115}+x_{115}^{2}+2x_{116}+x_{116}^{2}
\end{align*}
\begin{align*}
r_{102}&=
-x_{11}x_{20}-x_{12}x_{20}-x_{13}x_{20}-x_{14}x_{20}-x_{15}x_{20}-x_{16}x_{20}-x_{17}x_{20}-x_{18}x_{20} \\
&\quad -x_{19}x_{20}-10x_{20}+x_{118}
\end{align*}
\begin{align*}
r_{103}&=
1-x_{119}-x_{119}x_{164}-x_{119}x_{165}-x_{119}x_{215}+x_{170}+x_{171}+x_{177}+x_{182}+x_{187}+x_{218}
\end{align*}
\begin{align*}
r_{104}&=
-x_{20}-x_{20}x_{204}-x_{20}x_{205}-x_{20}x_{210}+x_{120}
\end{align*}
\begin{align*}
r_{105}&=
-80-2x_{0}-2x_{0}x_{123}-x_{1}-x_{1}x_{123}-x_{2}-x_{2}x_{123}-2x_{3}-2x_{3}x_{123}-x_{4}-x_{4}x_{123}-x_{5} \\
&\quad -x_{5}x_{123}-x_{6}-x_{6}x_{123}-x_{8}-x_{8}x_{123}-2x_{11}-2x_{11}x_{123}-14x_{12}-14x_{12}x_{123} \\
&\quad -x_{13}-x_{13}x_{123}-x_{14}-x_{14}x_{123}-x_{15}-x_{15}x_{123}-x_{16}-x_{16}x_{123}-x_{17}-x_{17}x_{123} \\
&\quad -x_{18}-x_{18}x_{123}-x_{19}-x_{19}x_{123}-x_{20}-x_{20}x_{123}-4x_{22}-4x_{22}x_{123}-2x_{25} \\
&\quad -2x_{25}x_{123}-x_{29}-x_{29}x_{123}-x_{31}-x_{31}x_{123}-x_{33}-x_{33}x_{123}-x_{35}-x_{35}x_{123} \\
&\quad -2x_{39}-2x_{39}x_{123}-4x_{41}-4x_{41}x_{123}-2x_{42}-2x_{42}x_{123}-x_{45}-x_{45}x_{123}-2x_{67} \\
&\quad -2x_{67}x_{123}-x_{118}-x_{118}x_{123}-x_{119}-x_{119}x_{123}-80x_{123}-x_{123}x_{164}-x_{123}x_{165} \\
&\quad -x_{123}x_{192}-x_{123}x_{193}-x_{123}x_{199}-x_{123}x_{204}-x_{123}x_{205}-x_{123}x_{210}-x_{123}x_{215} \\
&\quad +x_{124}-x_{164}-x_{165}-x_{192}-x_{193}-x_{199}-x_{204}-x_{205}-x_{210}-x_{215}
\end{align*}
\begin{align*}
r_{106}&=
-2x_{0}x_{124}-x_{1}x_{124}-x_{2}x_{124}-2x_{3}x_{124}-x_{4}x_{124}-x_{5}x_{124}-x_{6}x_{124}-x_{8}x_{124} \\
&\quad -2x_{11}x_{124}-14x_{12}x_{124}-x_{13}x_{124}-x_{14}x_{124}-x_{15}x_{124}-x_{16}x_{124}-x_{17}x_{124} \\
&\quad -x_{18}x_{124}-x_{19}x_{124}-x_{20}x_{124}-4x_{22}x_{124}-2x_{25}x_{124}-x_{29}x_{124}-x_{31}x_{124} \\
&\quad -x_{33}x_{124}-x_{35}x_{124}-2x_{39}x_{124}-4x_{41}x_{124}-2x_{42}x_{124}-x_{45}x_{124}-2x_{67}x_{124} \\
&\quad -x_{118}x_{124}-x_{119}x_{124}-81x_{124}-x_{124}x_{164}-x_{124}x_{165}-x_{124}x_{192}-x_{124}x_{193} \\
&\quad -x_{124}x_{199}-x_{124}x_{204}-x_{124}x_{205}-x_{124}x_{210}-x_{124}x_{215}+x_{125}
\end{align*}
\begin{align*}
r_{107}&=
-1+x_{121}^{2}+x_{124}^{2}-x_{125}^{2}
\end{align*}
\begin{align*}
r_{108}&=
-2-2x_{119}-x_{119}^{2}+2x_{119}x_{121}+2x_{119}x_{122}-x_{120}x_{126}+2x_{121}+2x_{122}
\end{align*}
\begin{align*}
r_{109}&=
-x_{129}^{2}+x_{220}
\end{align*}
\begin{align*}
r_{110}&=
1-x_{121}+x_{121}x_{220}-x_{122}+x_{122}x_{220}-4x_{128}x_{220}+x_{133}-x_{220}
\end{align*}
\begin{align*}
r_{111}&=
-x_{121}x_{128}-x_{122}x_{128}+x_{134}
\end{align*}
\begin{align*}
r_{112}&=
-1+x_{127}^{2}+x_{128}^{2}-x_{134}^{2}
\end{align*}
\begin{align*}
r_{113}&=
-x_{121}x_{130}-x_{122}x_{130}+x_{135}
\end{align*}
\begin{align*}
r_{114}&=
-1+x_{129}^{2}+x_{130}^{2}-x_{135}^{2}
\end{align*}
\begin{align*}
r_{115}&=
-x_{132}-x_{132}x_{133}+x_{136}
\end{align*}
\begin{align*}
r_{116}&=
-1+x_{131}^{2}+x_{132}^{2}-x_{136}^{2}
\end{align*}
\begin{align*}
r_{117}&=
-x_{128}^{2}+x_{138}
\end{align*}
\begin{align*}
r_{118}&=
x_{130}-2x_{137}x_{138}-2x_{138}
\end{align*}
\begin{align*}
r_{119}&=
-x_{127}-x_{129}x_{139}+x_{131}
\end{align*}
\begin{align*}
r_{120}&=
-25-x_{11}-13x_{12}-4x_{22}-2x_{25}-x_{29}-x_{31}-x_{33}-x_{35}-2x_{39}-4x_{41}-2x_{42}-x_{45}-2x_{67} \\
&\quad -4x_{128}x_{140}+x_{132}
\end{align*}
\begin{align*}
r_{121}&=
-x_{141}x_{142}+x_{143}
\end{align*}
\begin{align*}
r_{122}&=
-1-2x_{121}x_{143}-2x_{122}x_{143}+x_{141}^{2}+x_{142}^{2}
\end{align*}
\begin{align*}
r_{123}&=
-2-2x_{3}-x_{3}^{2}+2x_{3}x_{121}+2x_{3}x_{122}+2x_{121}+2x_{122}-x_{144}
\end{align*}
\begin{align*}
r_{124}&=
x_{3}x_{141}-x_{4}+x_{141}-x_{142}-x_{144}x_{145}
\end{align*}
\begin{align*}
r_{125}&=
-x_{146}x_{147}+x_{148}
\end{align*}
\begin{align*}
r_{126}&=
-1-2x_{121}x_{148}-2x_{122}x_{148}+x_{146}^{2}+x_{147}^{2}
\end{align*}
\begin{align*}
r_{127}&=
-2-2x_{5}-x_{5}^{2}-2x_{5}x_{6}+2x_{5}x_{121}+2x_{5}x_{122}-2x_{6}-x_{6}^{2}+2x_{6}x_{121}+2x_{6}x_{122} \\
&\quad +2x_{121}+2x_{122}-x_{149}
\end{align*}
\begin{align*}
r_{128}&=
x_{5}x_{146}+x_{6}x_{146}-x_{20}+x_{146}-x_{147}-x_{149}x_{150}
\end{align*}
\begin{align*}
r_{129}&=
-1-x_{118}^{2}+2x_{118}x_{121}+2x_{118}x_{122}-x_{151}
\end{align*}
\begin{align*}
r_{130}&=
x_{118}x_{146}-x_{119}-x_{147}-x_{151}x_{152}
\end{align*}
\begin{align*}
r_{131}&=
-1-x_{119}^{2}+2x_{119}x_{121}+2x_{119}x_{122}-x_{153}
\end{align*}
\begin{align*}
r_{132}&=
-1+x_{119}x_{141}-x_{142}-x_{153}x_{154}-x_{164}-x_{165}-x_{215}
\end{align*}
\begin{align*}
r_{133}&=
-x_{155}x_{156}+x_{157}
\end{align*}
\begin{align*}
r_{134}&=
-1-2x_{121}x_{157}-2x_{122}x_{157}+x_{155}^{2}+x_{156}^{2}
\end{align*}
\begin{align*}
r_{135}&=
-1+x_{119}x_{155}-x_{153}x_{158}-x_{156}-x_{192}-x_{193}-x_{199}
\end{align*}
\begin{align*}
r_{136}&=
-1+x_{119}x_{146}-x_{147}-x_{153}x_{159}-x_{204}-x_{205}-x_{210}
\end{align*}
\begin{align*}
r_{137}&=
-x_{160}x_{161}+x_{162}
\end{align*}
\begin{align*}
r_{138}&=
-1-2x_{121}x_{162}-2x_{122}x_{162}+x_{160}^{2}+x_{161}^{2}
\end{align*}
\begin{align*}
r_{139}&=
-1-x_{5}-x_{5}x_{164}-x_{5}x_{165}-x_{5}x_{215}-x_{6}-x_{6}x_{164}-x_{6}x_{165}-x_{6}x_{215}+x_{119}x_{160} \\
&\quad +x_{160}-x_{161}-x_{163}-x_{163}x_{170}-x_{163}x_{171}-x_{163}x_{177}-x_{163}x_{182}-x_{163}x_{187} \\
&\quad -x_{163}x_{218}-2x_{164}-x_{165}-x_{215}
\end{align*}
\begin{align*}
r_{140}&=
-x_{166}x_{167}+x_{168}
\end{align*}
\begin{align*}
r_{141}&=
-1-2x_{121}x_{168}-2x_{122}x_{168}+x_{166}^{2}+x_{167}^{2}
\end{align*}
\begin{align*}
r_{142}&=
-x_{20}-x_{20}x_{170}-x_{20}x_{171}-x_{20}x_{177}-x_{20}x_{182}-x_{20}x_{187}-x_{20}x_{218}+x_{172}
\end{align*}
\begin{align*}
r_{143}&=
x_{119}x_{166}+x_{166}-x_{167}-x_{169}x_{172}-x_{170}
\end{align*}
\begin{align*}
r_{144}&=
-x_{173}x_{174}+x_{175}
\end{align*}
\begin{align*}
r_{145}&=
-1-2x_{121}x_{175}-2x_{122}x_{175}+x_{173}^{2}+x_{174}^{2}
\end{align*}
\begin{align*}
r_{146}&=
x_{119}x_{173}-x_{172}x_{176}+x_{173}-x_{174}-x_{177}
\end{align*}
\begin{align*}
r_{147}&=
-x_{178}x_{179}+x_{180}
\end{align*}
\begin{align*}
r_{148}&=
-1-2x_{121}x_{180}-2x_{122}x_{180}+x_{178}^{2}+x_{179}^{2}
\end{align*}
\begin{align*}
r_{149}&=
x_{119}x_{178}-x_{172}x_{181}+x_{178}-x_{179}-x_{182}
\end{align*}
\begin{align*}
r_{150}&=
-x_{183}x_{184}+x_{185}
\end{align*}
\begin{align*}
r_{151}&=
-1-2x_{121}x_{185}-2x_{122}x_{185}+x_{183}^{2}+x_{184}^{2}
\end{align*}
\begin{align*}
r_{152}&=
x_{119}x_{183}-x_{172}x_{186}+x_{183}-x_{184}-x_{187}
\end{align*}
\begin{align*}
r_{153}&=
-x_{188}x_{189}+x_{190}
\end{align*}
\begin{align*}
r_{154}&=
-1-2x_{121}x_{190}-2x_{122}x_{190}+x_{188}^{2}+x_{189}^{2}
\end{align*}
\begin{align*}
r_{155}&=
-x_{20}-x_{20}x_{192}-x_{20}x_{193}-x_{20}x_{199}+x_{194}
\end{align*}
\begin{align*}
r_{156}&=
x_{119}x_{188}+x_{188}-x_{189}-x_{191}x_{194}-x_{192}
\end{align*}
\begin{align*}
r_{157}&=
-x_{195}x_{196}+x_{197}
\end{align*}
\begin{align*}
r_{158}&=
-1-2x_{121}x_{197}-2x_{122}x_{197}+x_{195}^{2}+x_{196}^{2}
\end{align*}
\begin{align*}
r_{159}&=
x_{119}x_{195}-x_{194}x_{198}+x_{195}-x_{196}-x_{199}
\end{align*}
\begin{align*}
r_{160}&=
-x_{200}x_{201}+x_{202}
\end{align*}
\begin{align*}
r_{161}&=
-1-2x_{121}x_{202}-2x_{122}x_{202}+x_{200}^{2}+x_{201}^{2}
\end{align*}
\begin{align*}
r_{162}&=
x_{119}x_{200}-x_{120}x_{203}+x_{200}-x_{201}-x_{204}
\end{align*}
\begin{align*}
r_{163}&=
-x_{206}x_{207}+x_{208}
\end{align*}
\begin{align*}
r_{164}&=
-1-2x_{121}x_{208}-2x_{122}x_{208}+x_{206}^{2}+x_{207}^{2}
\end{align*}
\begin{align*}
r_{165}&=
x_{119}x_{206}-x_{120}x_{209}+x_{206}-x_{207}-x_{210}
\end{align*}
\begin{align*}
r_{166}&=
-x_{211}x_{212}+x_{213}
\end{align*}
\begin{align*}
r_{167}&=
-1-2x_{121}x_{213}-2x_{122}x_{213}+x_{211}^{2}+x_{212}^{2}
\end{align*}
\begin{align*}
r_{168}&=
-x_{20}-x_{20}x_{164}-x_{20}x_{165}-x_{20}x_{215}+x_{216}
\end{align*}
\begin{align*}
r_{169}&=
x_{119}x_{211}+x_{211}-x_{212}-x_{214}x_{216}-x_{215}
\end{align*}
\begin{align*}
r_{170}&=
x_{119}x_{128}+x_{127}+x_{128}-x_{134}-x_{172}x_{217}-x_{218}
\end{align*}
\begin{align*}
r_{171}&=
-326-48x_{0}-2x_{0}^{2}-2x_{0}x_{3}+2x_{0}x_{141}+2x_{0}x_{160}-x_{1}^{2}+2x_{1}x_{146}-4x_{2}-x_{2}^{2} \\
&\quad -2x_{2}x_{8}+2x_{2}x_{155}-24x_{3}-x_{3}^{2}+2x_{3}x_{160}-4x_{8}-x_{8}^{2}+2x_{8}x_{155}-3x_{11} \\
&\quad -x_{11}^{2}+2x_{11}x_{166}-15x_{12}-x_{12}^{2}+2x_{12}x_{173}-2x_{13}-x_{13}^{2}+2x_{13}x_{178}-2x_{14} \\
&\quad -x_{14}^{2}+2x_{14}x_{183}-2x_{15}-x_{15}^{2}+2x_{15}x_{188}-2x_{16}-x_{16}^{2}+2x_{16}x_{195}-2x_{17} \\
&\quad -x_{17}^{2}+2x_{17}x_{200}-2x_{18}-x_{18}^{2}+2x_{18}x_{206}-2x_{19}-x_{19}^{2}+2x_{19}x_{211}-4x_{22} \\
&\quad -2x_{25}-x_{29}-x_{31}-x_{33}-x_{35}-2x_{39}-4x_{41}-2x_{42}-x_{45}-2x_{67}-x_{121}x_{219}-x_{122}x_{219} \\
&\quad +x_{128}+24x_{141}-x_{141}^{2}-x_{146}^{2}+4x_{155}-x_{155}^{2}+24x_{160}-x_{160}^{2}+2x_{166} \\
&\quad -x_{166}^{2}+2x_{173}-x_{173}^{2}+2x_{178}-x_{178}^{2}+2x_{183}-x_{183}^{2}+2x_{188}-x_{188}^{2}+2x_{195} \\
&\quad -x_{195}^{2}+2x_{200}-x_{200}^{2}+2x_{206}-x_{206}^{2}+2x_{211}-x_{211}^{2}+x_{219}
\end{align*}

\endgroup

\clearpage
\begin{thebibliography}{9}
\bibitem{repo}
V. Reshetnikov, \emph{ProveIt: DiophantineConsistency}, pinned revision
\texttt{e0c98c629bf10}, accessed 25 October 2026.
\url{https://github.com/VladimirReshetnikov/ProveIt/tree/e0c98c629bf10d30b695884f6b1d2ab361bdf513/Computability/DiophantineConsistency}.

\bibitem{incoming}
ProveIt incoming research archives, same pinned revision. In particular,
\emph{ConZFC Certificate Reductions} (25 October),
\emph{ConZFC Cofinal Anchors, Height One, Round 17} (25 October), and
\emph{conzfc graph digits} (24 October). Exact archive and inherited-file
hashes are retained in the component provenance records.
\url{https://github.com/VladimirReshetnikov/ProveIt/tree/e0c98c629bf10d30b695884f6b1d2ab361bdf513/docs/incoming}.

\bibitem{p66}
ProveIt, \emph{A 66-operation consistency polynomial by a squared-modulus
quotient}, package 51; and \emph{one-page-compact66}, the complete indexed
specialization. Both are read at the same pinned revision; the latter's
full ordinary proof and polynomial are retained in this archive.
\url{https://github.com/VladimirReshetnikov/ProveIt/tree/e0c98c629bf10d30b695884f6b1d2ab361bdf513/Computability/DiophantineConsistency/candidates/one-page-compact66}.

\bibitem{jones}
J. P. Jones, Universal Diophantine equation, \emph{Journal of Symbolic Logic}
47(3) (1982), 549--571. \url{https://doi.org/10.2307/2273588}.

\bibitem{bayer}
J. Bayer, M. David, M. Hassler, Y. Matiyasevich, and D. Schleicher,
\emph{Diophantine Equations over $\Z$: Universal Bounds and Parallel
Formalization}, arXiv:2506.20909, version 3, 22 October 2026.
\url{https://arxiv.org/abs/2506.20909v3}.
\end{thebibliography}
\end{document}
